\section{The checked compiler spine}

\subsection{Simultaneous discharge of generated assumptions}

The central structural lemma is a finite list form of cut.  It is almost a
direct call to the checked theorem \path{KPrf.cut}.

\begin{theorem}[List cut]\label{thm:list-cut}
Let \(\Gamma\) be a base clause context, let
\(\Delta=[B_0,\ldots,B_{r-1}]\) be native formulas, and let \(A\) be a native
formula.  If
\[
  h:\KPrf(\Delta^{\mathsf{asm}}\mathbin{++}\Gamma,A)
\]
and for every \(B_i\in\Delta\) there is
\(h_i:\KPrf(\Gamma,B_i)\), then \(\KPrf(\Gamma,A)\).
\end{theorem}

\begin{proof}
Apply \path{KPrf.cut} to \(h\) with target context \(\Gamma\).  Every source
clause belongs either to \(\Gamma\), in which case it is retained, or to
\(\Delta^{\mathsf{asm}}\), in which case it has the form
\(\Cl.\mathsf{asm}(B_i)\) and is discharged by \(h_i\).  These are exactly the
two alternatives required by the checked cut theorem.  The argument is a
finite list-membership decomposition and does not inspect the internal
structure of \(h\).
\end{proof}

A candidate Lean proof is included in \path{lean/NativeClauseFramework/Bridge.lean}:

\begin{lstlisting}[language={},caption={Candidate direct proof of list cut}]
theorem cutAsmList
    (h : KPrf (Delta.map Cl.asm ++ Gamma) A)
    (hs : forall X, X in Delta -> KPrf Gamma X) : KPrf Gamma A := by
  apply h.cut
  intro c hc
  rw [List.mem_append] at hc
  rcases hc with hc | hc
  case left => obtain <X, hX, rfl> := List.mem_map.1 hc
    exact .inr <X, rfl, hs X hX>
  case right => exact .inl hc
\end{lstlisting}
The archive uses ASCII approximations in explanatory snippets where convenient;
the actual Lean file uses Lean syntax.

\subsection{Acyclic program compilation}

A program is given in topological order.  Node \(i\) may depend on any subset
of outputs \(A_0,\ldots,A_{i-1}\), and its local proof is allowed to assume all
previous outputs for simplicity:
\[
 h_i:\KPrf([A_0,\ldots,A_{i-1}]^{\mathsf{asm}}
                 \mathbin{++}\Gamma,A_i).
\]
Unused assumptions are harmless.

\begin{theorem}[Native DAG compilation]\label{thm:dag-compile}
For every topologically ordered finite list of native clause nodes over
\(\Gamma\), the compiler constructs a derivation
\(\KPrf(\Gamma,A_i)\) for every node output.  In particular it constructs a
derivation of the designated terminal output.
\end{theorem}

\begin{proof}
Induct on the node list.  The empty list is immediate.  At node \(i\), the
induction hypothesis supplies derivations of every earlier \(A_j\) from
\(\Gamma\).  Apply Theorem~\ref{thm:list-cut} to the node's local derivation
\(h_i\).  This yields \(\KPrf(\Gamma,A_i)\), which is appended to the compiled
proof environment.  No semantic property of the executable node is used.
\end{proof}

The theorem remains true when a node lists only its actual predecessors: map
those predecessors into a temporary assumption list and apply list cut.  A
reverse index or membership proof is metadata used to route the corresponding
compiled derivations.

\subsection{Specialization and parameter families}

A local theorem is normally proved with free variables and reused at many
concrete terms.  The checked theorem \path{KPrf.subst} states that substitution
of terms for all free variables transports both the clause context and the
conclusion.

\begin{theorem}[Specialization]
If \(h:\KPrf(\Gamma,A)\) and
\(\sigma:\mathrm{Fin}(n)\to\mathrm{Tm}(m)\), then
\[
 \KPrf(\Gamma[\sigma],A[\sigma]).
\]
\end{theorem}

\begin{proof}
This is exactly \path{KPrf.subst}.  Its checked proof is by induction on the
five rule families and includes the nontrivial binder bookkeeping for
elimination and introduction.
\end{proof}

This theorem is the formal reason to compile a library of generic node lemmas.
For example, a generic rational-addition contract with variables for four
numerators and denominators can be specialized to every occurrence in a graph
without redoing its proof.

\subsection{Export to ordinary iPRA}

The checked method \path{KPrf.toNPrf} maps a clause derivation to a native
scalar-core derivation whose assumptions are the ordinary native readings of
the clauses.  The theorem \path{native_iff_pra} maps that derivation to ordinary
iPRA.

\begin{theorem}[Clause-to-iPRA export]\label{thm:export}
For \(h:\KPrf(\Gamma,A)\), there is a derivation
\[
  \Gamma.\mathsf{map}(c\mapsto (c.\mathsf{toN}).\mathsf{toPRA})
  \ \vdash_{\PRAAx}\ A.\mathsf{toPRA}.
\]
In particular, a closed \(\KPrf([],A)\) yields a closed ordinary iPRA theorem.
\end{theorem}

\begin{proof}
Apply \path{KPrf.toNPrf} to obtain
\(\NPrf(\Gamma.\mathsf{map}(\Cl.\mathsf{toN}),A)\).  Apply the forward
direction of \path{native_iff_pra}.  The resulting assumption list is the
stated list after associativity of \path{List.map}.
\end{proof}

This theorem gives the complete trusted route
\[
  \text{local }\KPrf
   \xrightarrow{\text{cut/subst}}
  \text{terminal }\KPrf
   \xrightarrow{\mathsf{toNPrf}}
  \NPrf
   \xrightarrow{\mathsf{native\_iff\_pra}}
  \Prf\;\PRAAx.
\]

\subsection{Import of direct iPRA derivations}

The route also runs in the opposite direction.  The theorem
\path{pra_iff_native} converts an ordinary derivation to a native derivation,
and \path{kprf_iff_nprf} converts the native derivation to the clause calculus
after \path{unQ} normalization.

\begin{theorem}[iPRA-to-clause import]\label{thm:import}
If \(h:\Gamma\vdash_{\PRAAx}\varphi\), then
\[
 \KPrf\left(
  (\Gamma.\mathsf{map}\,\NFm.\mathsf{ofPRA}).
      \mathsf{map}(X\mapsto\Cl.\mathsf{asm}(\NFm.\mathsf{unQ}(X))),
  \NFm.\mathsf{unQ}(\NFm.\mathsf{ofPRA}(\varphi))\right).
\]
\end{theorem}

\begin{proof}
The forward direction of \path{pra_iff_native} gives
\(\NPrf(\Gamma.\mathsf{map}\,\NFm.\mathsf{ofPRA},
\NFm.\mathsf{ofPRA}(\varphi))\).  Apply the reverse direction of
\path{kprf_iff_nprf}.
\end{proof}

Thus a theorem produced by \path{ipra_ring} may be inserted as a local native
node, used by search or branching code, and exported again after composition.
This is the main bridge between the new direct derivation library and the
finite-control framework.

\subsection{Compiler correctness is syntactic}

The compiler theorem is stronger than a semantic soundness statement in one
respect and weaker in another.

It is stronger because the output is an actual derivation, not merely a proof
that every model satisfying the inputs satisfies the output.  It is weaker
because the compiler assumes local derivations rather than inferring them from
arbitrary successful computations.  That is the correct boundary.  A small
trusted proof calculus should not contain a general rule ``verified by external
program.''

\subsection{Demand compilation and proof compilation are separate}

The host compiler performs two passes.

\begin{enumerate}
\item The \emph{demand pass} computes the finite precision, fuel, and resource
requirements needed by each node.  Its correctness is a theorem about finite
natural-number algorithms.
\item The \emph{proof pass} instantiates local node theorems at the computed
terms and folds them with list cut.
\end{enumerate}

The separation prevents circularity.  The demand pass may be executed in the
host to generate concrete numerals, but if the theorem quantifies over the
requested precision, the demand function itself is represented by a PR symbol
and its defining equations are used in the object proof.  In either case, the
proof pass never assumes the semantic correctness of an opaque host result.

\subsection{No need to encode the graph as a number}

The general theorem is implemented at the Lean metalevel over a finite list of
nodes.  For a particular application, Lean elaborates the list and constructs a
proof term.  The object theory sees only the resulting formulas, PR terms, and
proof.  Encoding the graph and writing an internal graph interpreter would add
substantial machinery and would approach the forbidden universal-evaluator
problem without providing additional mathematical strength.

This staged approach is analogous to a verified macro expansion: generic Lean
code produces a concrete object-language derivation, and the small inductive
proof kernel checks the result.
